\chapter{Vertering en voedingsstoffen}\label{ch:g7:digestion-and-nutrients}

\Cref{ch:g4:journey-of-food} volgde de hap brood door de buis omlaag en
beloofde dat ``\omterm{def:g4:journey-of-food:digestion}{verteringssappen} hem uit elkaar halen''. Dit jaar zetten we
de deuren van de werkplaats open: wat doen die sappen precies, wat zijn de
\omterm{def:g7:digestion-and-nutrients:families}{voedingsstoffen} die ze opleveren, en hoe klaart de wand van de \omterm{def:g4:journey-of-food:tract}{dunne darm}
de grootste grensovergang van het lichaam? De hap brood is terug --- onder
betere instrumenten.

\section{De sappen aan het werk}

\begin{proposition}[Vertering is chemisch uit elkaar halen]\label{prop:g7:digestion-and-nutrients:chemical}
Kauwen en kneden maken voedsel alleen \emph{kleiner}; het echte uit elkaar
halen is chemisch. De \emph{verteringssappen}\index{verteringssap} ---
\omterm{def:g4:journey-of-food:tract}{speeksel}, het maagsap, de sappen die in de \omterm{def:g4:journey-of-food:tract}{dunne darm} gegoten worden ---
bevatten werkzame stoffen die de grote bestanddelen van voedsel tot kleine
knippen: zetmeel tot suikers, de bouwstof van vlees en eieren tot haar
kleine eenheden, vetten tot de hunne. Elk sap heeft zijn specialiteiten,
en elk werkt het best bij de \omterm{def:g6:exploring-our-environment:conditions}{omstandigheden} van zijn eigen station.
\end{proposition}
\admitted

\begin{example}[De versie op de werkbank]\label{ex:g7:digestion-and-nutrients:bench}
De klassieke proef laat de \omterm{def:g4:journey-of-food:digestion}{vertering} in glas draaien. Twee buizen met
gekookte zetmeelpap op lichaamstemperatuur: de een krijgt vers \omterm{def:g4:journey-of-food:tract}{speeksel},
de ander gekookt (bedorven) \omterm{def:g4:journey-of-food:tract}{speeksel}. Een uur later vindt de zetmeeltest
buis 2 onveranderd --- nog steeds zetmeel --- terwijl de pap in buis 1 in
suikers veranderd is, precies de zoetheid uit
\cref{ex:g4:journey-of-food:sweet}, zonder ook maar één \omterm{def:g4:journey-of-food:tract}{mond}
voortgebracht. Het sap alleen deed het uit elkaar halen --- en de gekookte
controle laat zien dat de werkzame stof ervan door hitte vernield wordt,
zoals fijne machinerie dat zou worden.
\end{example}

\begin{remark}[Waarom eigenlijk uit elkaar halen?]\label{rem:g7:digestion-and-nutrients:why}
De grote bestanddelen van voedsel kunnen de darmwand net zomin passeren
als een kleerkast een brievenbus. En het lichaam wil ze niet in hun
geheel: het wil \emph{onderdelen} --- kleine eenheden waarmee het zijn
eigen stof kan opbouwen
(\cref{prop:g6:matter-from-food:transform}). \omterm{def:g4:journey-of-food:digestion}{Vertering} is uit elkaar halen
voor het vervoer \emph{en} voor de opbouw: meubels tot bouwpakket, bij
elke maaltijd.
\end{remark}

\section{De voedingsstoffen}

\begin{definition}[De families van voedingsstoffen]\label{def:g7:digestion-and-nutrients:families}
Het uit elkaar halen levert de
\emph{voedingsstoffen}\index{voedingsstof} op --- de kleine bestanddelen
die klaar zijn voor het \omterm{prop:g4:journey-of-food:absorb}{bloed}:
\begin{itemize}
  \item \emph{suikers} (vooral \omterm{prop:g7:organs-at-work:bill}{glucose}, de brandstof uit
        \cref{prop:g7:organs-at-work:bill}), uit zetmeelrijk en zoet
        voedsel;
  \item de kleine eenheden van de \emph{bouwstof}, uit vlees, vis,
        eieren, zuivel en peulvruchten --- de bakstenen van de bouw uit
        \cref{prop:g3:balanced-diet:jobs};
  \item de eenheden van \emph{vetten} --- geconcentreerde brandstof en
        bouwmateriaal tegelijk;
  \item plus passagiers die geen uiteenhalen nodig hebben:
        \emph{vitamines}, \emph{\omterm{prop:g4:what-plants-need:needs}{minerale zouten}} en het water zelf.
\end{itemize}
De \omterm{def:g3:balanced-diet:families}{voedingsgroepen} uit \cref{def:g3:balanced-diet:families} zijn in de
grond leveranciers van deze families van voedingsstoffen --- het
boodschappenlijstje achter het boodschappenlijstje.
\end{definition}

\begin{example}[Eén lunch, uitgesplitst]\label{ex:g7:digestion-and-nutrients:lunch}
De evenwichtige lunch uit \cref{ex:g3:balanced-diet:lunch}, opnieuw
gelezen als \omterm{def:g7:digestion-and-nutrients:families}{voedingsstoffen}: pasta --- zetmeel, bij de liter tot \omterm{prop:g7:organs-at-work:bill}{glucose}
uit elkaar gehaald; vis --- bouwstof, tot haar eenheden geknipt; de olie
op de bonen --- vetteenheden; bonen en appel --- vitamines, mineralen,
water, plus vezels. En let op de vezels: die zijn de uitzondering die het
systeem in kaart brengt --- geen enkel sap van ons kan ze knippen, dus
steken ze niets over en reizen ze door --- de massa van de restanten
(\cref{prop:g4:journey-of-food:absorb}), die haar eigen nuttige veegwerk
doet.
\end{example}

\section{De grote oversteek}

\begin{proposition}[Opname bij de darmvlokken]\label{prop:g7:digestion-and-nutrients:absorption}
De oversteek die in \cref{prop:g4:journey-of-food:absorb} beloofd werd,
gebeurt door de binnenwand van de \omterm{def:g4:journey-of-food:tract}{dunne darm}, en die wand is ervoor
gebouwd: geplooid, en bekleed met miljoenen
\emph{darmvlokken}\index{darmvlok} --- vingervormige bobbeltjes van een
millimeter hoog, elk met zijn eigen fijne \omterm{def:g4:heart-and-blood:vessels}{bloedvaten} erin. De \omterm{prop:g7:digestion-and-nutrients:absorption}{darmvlokken}
zijn de reden dat het oversteekoppervlak de omvang van een kamertapijt
haalt (\cref{rem:g4:journey-of-food:surface}): plooien op plooien op
vingers. \omterm{def:g7:digestion-and-nutrients:families}{Voedingsstoffen} steken hun dunne wanden over naar het \omterm{prop:g4:journey-of-food:absorb}{bloed}, en
het \omterm{prop:g4:journey-of-food:absorb}{bloed} brengt ze eerst naar een groot sorteer- en opslagorgaan --- de
\emph{lever}\index{lever} --- voordat ze algemeen bezorgd worden.
\end{proposition}
\admitted

\begin{omfigure}
\begin{tikzpicture}[scale=0.9]
  % intestine wall with villi
  \draw[thick] (0,0) -- (0,3.4);
  \draw[thick] (0,0) -- (8.2,0);
  \node[font=\small, rotate=90] at (-0.4,1.7) {darmwand};
  % villi
  \foreach \x in {0.9,2.4,3.9,5.4,6.9} {
    \draw[thick, fill=omDef!8] (\x,0)
      .. controls ({\x-0.42},1.4) and ({\x-0.42},2.4) .. (\x,2.6)
      .. controls ({\x+0.42},2.4) and ({\x+0.42},1.4) .. (\x,0);
    % vessels inside
    \draw[omThm, thick] ({\x-0.14},0.15)
      .. controls ({\x-0.18},1.3) .. (\x,2.25)
      .. controls ({\x+0.18},1.3) .. ({\x+0.14},0.15);
  }
  \node[font=\footnotesize, align=center] at (4.1,-0.6)
    {elke darmvlok bevat zijn eigen fijne bloedvaten};
  % nutrients crossing
  \foreach \p in {(1.7,2.9),(3.2,3.0),(4.7,2.85),(6.2,2.95)} {
    \fill[omMeth] \p circle (0.07);
  }
  \node[font=\footnotesize, omMeth] at (6.9,3.25)
    {voedingsstoffen in de buis};
  \draw[->, omMeth, thick] (3.2,2.9) -- (3.65,2.2);
  \node[font=\footnotesize, align=left] at (10.0,1.7)
    {de oversteek:\\ door de dunne wand\\ van de darmvlok,\\ het bloed in};
\end{tikzpicture}

{\small Het binnentapijt van de darm: miljoenen \omterm{prop:g7:digestion-and-nutrients:absorption}{darmvlokken}, elk een
dunwandige vinger vol \omterm{def:g4:heart-and-blood:vessels}{bloedvaten}. Oppervlak op oppervlak gevouwen --- de
grote grensovergang van het lichaam.}
\end{omfigure}

\begin{example}[Opnieuw het bestek]\label{ex:g7:digestion-and-nutrients:spec}
Groot, dun, vochtig, met \omterm{prop:g4:journey-of-food:absorb}{bloed} doorregen: de \omterm{prop:g7:digestion-and-nutrients:absorption}{darmvlokken} voldoen precies
aan het bestek van \cref{prop:g7:how-animals-breathe:spec} --- voor
\omterm{def:g7:digestion-and-nutrients:families}{voedingsstoffen} in plaats van \omterm{def:g7:what-is-respiration:respiration}{zuurstof}. De zakjes van de \omterm{def:g7:how-animals-breathe:four}{longen}, de \omterm{def:g1:animals-around-us:covering}{veren}
van de \omterm{def:g7:how-animals-breathe:four}{kieuwen}, de vlokken van de darm: waar het lichaam veel over een
grens moet brengen, bouwt het dezelfde machine. Herken het bestek één
keer, en je herkent het overal.
\end{example}

\begin{method}[Een voedingsstof van begin tot eind volgen]\label{met:g7:digestion-and-nutrients:trace}
Volg bij elke hap één \omterm{def:g7:digestion-and-nutrients:families}{voedingsstof}:
\begin{enumerate}
  \item noem de familie van het voedsel en zijn hoofdbestanddelen;
  \item noteer station voor station waar het uit elkaar halen begint en
        eindigt, en welke sappen eraan gewerkt hebben;
  \item steek er bij de \omterm{prop:g7:digestion-and-nutrients:absorption}{darmvlokken} mee over, het \omterm{prop:g4:journey-of-food:absorb}{bloed} in, via de \omterm{prop:g7:digestion-and-nutrients:absorption}{lever};
  \item bezorg hem: bij een werkende \omterm{def:g5:first-look-at-cells:cell}{cel} als brandstof
        (\cref{prop:g7:organs-at-work:bill}), of bij een bouwplaats als
        materiaal (\cref{def:g6:matter-from-food:production}).
\end{enumerate}
\end{method}

\begin{example}[Het spoor, gevolgd op brood]\label{ex:g7:digestion-and-nutrients:breadtrace}
Brood: de zetmeelfamilie, vooral zetmeel. Het uit elkaar halen begint in
de \omterm{def:g4:journey-of-food:tract}{mond} (\omterm{def:g4:journey-of-food:tract}{speeksel} --- de zoetheid van lang kauwen), pauzeert in het zuur
van de \omterm{def:g4:journey-of-food:tract}{maag} en wordt afgemaakt in de sappen van de \omterm{def:g4:journey-of-food:tract}{dunne darm}: \omterm{prop:g7:organs-at-work:bill}{glucose}.
Oversteek: \omterm{prop:g7:digestion-and-nutrients:absorption}{darmvlokken}, \omterm{prop:g4:journey-of-food:absorb}{bloed}, \omterm{prop:g7:digestion-and-nutrients:absorption}{lever} --- die een deel van de \omterm{prop:g7:organs-at-work:bill}{glucose} voor
later op de bank zet, een van haar vele diensten. Bezorging: een dijspier
midden in een sprint verbrandt hem binnen het uur. De hap uit leerjaar
vier, volledig verantwoord.
\end{example}

\section{Oefeningen}

\begin{exercise}[$\star$]\label{exo:g7:digestion-and-nutrients:1}
Wat doen \omterm{def:g4:journey-of-food:digestion}{verteringssappen} wat tanden en kneden niet kunnen?
\end{exercise}

\begin{exercise}[$\star$]\label{exo:g7:digestion-and-nutrients:2}
Beschrijf de speekselproef met twee buizen: opzet, uitkomst en allebei de
conclusies.
\end{exercise}

\begin{exercise}[$\star$]\label{exo:g7:digestion-and-nutrients:3}
Noem de families van \omterm{def:g7:digestion-and-nutrients:families}{voedingsstoffen} en van elk één voedselbron.
\end{exercise}

\begin{exercise}[$\star$]\label{exo:g7:digestion-and-nutrients:4}
Welke bestanddelen van voedsel hoeven niet uit elkaar gehaald te worden?
En welk bestanddeel verslaat onze sappen volledig --- met welk nut?
\end{exercise}

\begin{exercise}[$\star$]\label{exo:g7:digestion-and-nutrients:5}
Wat is een darmvlok? Wat zit er in elke vlok?
\end{exercise}

\begin{exercise}[$\star$]\label{exo:g7:digestion-and-nutrients:6}
Waar gaat het \omterm{prop:g4:journey-of-food:absorb}{bloed} uit de darm heen vóór de algemene bezorging, en noem
één dienst die dat orgaan verleent.
\end{exercise}

\begin{exercise}[$\star\star$]\label{exo:g7:digestion-and-nutrients:7}
Leg het beeld van het bouwpakket uit
\cref{rem:g7:digestion-and-nutrients:why} uit: de twee redenen waarom
voedsel uit elkaar gehaald moet worden.
\end{exercise}

\begin{exercise}[$\star\star$]\label{exo:g7:digestion-and-nutrients:8}
Waarom doet de buis met gekookt \omterm{def:g4:journey-of-food:tract}{speeksel} ertoe? Noem de methode, en wat
het koken laat zien over de werkzame stof van het sap.
\end{exercise}

\begin{exercise}[$\star\star$]\label{exo:g7:digestion-and-nutrients:9}
Zet de \omterm{prop:g7:digestion-and-nutrients:absorption}{darmvlokken} naast de zakjes van de \omterm{def:g7:how-animals-breathe:four}{longen} en de \omterm{def:g1:animals-around-us:covering}{veren} van de
\omterm{def:g7:how-animals-breathe:four}{kieuwen}: noem het gedeelde bestek, en de vracht van elke machine.
\end{exercise}

\begin{exercise}[$\star\star$]\label{exo:g7:digestion-and-nutrients:10}
Laat \cref{met:g7:digestion-and-nutrients:trace} lopen over een gekookt \omterm{def:g2:animals-grow-up:born}{ei}
(bouwstof), tot aan een herstelplek in een genezende schaafwond
(\cref{exo:g6:cell-unit-of-life:10}).
\end{exercise}

\begin{exercise}[$\star\star$]\label{exo:g7:digestion-and-nutrients:11}
Een crashdieet slaat wekenlang alle bouwstofrijke voeding over. Voorspel
met \cref{def:g7:digestion-and-nutrients:families} en
\cref{def:g6:matter-from-food:production} wat er tekort raakt en wat er
het eerst onder lijdt in een lichaam dat nog groeit.
\end{exercise}

\begin{exercise}[$\star\star\star$]\label{exo:g7:digestion-and-nutrients:12}
``De \omterm{def:g4:journey-of-food:digestion}{vertering} maakt van een maaltijd een bezorglijst.'' Verdedig die zin
van begin tot eind --- sappen, \omterm{def:g7:digestion-and-nutrients:families}{voedingsstoffen}, \omterm{prop:g7:digestion-and-nutrients:absorption}{darmvlokken}, \omterm{prop:g7:digestion-and-nutrients:absorption}{lever},
bezorgingen --- in één alinea.
\end{exercise}

\section{Probleem: De glazen darm}

\begin{problem}[{Weekendprobleem --- de vertering op de werkbank gebouwd,
buis voor buis}]\label{pb:g7:digestion-and-nutrients:1}
De klas bouwt een ``glazen darm'': een rek met buizen op
lichaamstemperatuur, die elk één station van het kanaal nabootsen. Buis M
(\omterm{def:g4:journey-of-food:tract}{mond}): zetmeelpap plus vers \omterm{def:g4:journey-of-food:tract}{speeksel}. Buis M$'$: zetmeelpap plus gekookt
\omterm{def:g4:journey-of-food:tract}{speeksel}. Buis S (\omterm{def:g4:journey-of-food:tract}{maag}): blokjes eiwit plus maagsap van de
apotheek. Buis S$'$: dezelfde blokjes in gewoon water. Alle buizen draaien
twee uur; daarna de tests --- op zetmeel, op suikers, en een blik op de
blokjes.

\medskip\noindent\textbf{Deel I --- Voorspellingen en
uitkomsten.}
\begin{enumerate}
  \item Voorspel de uitkomsten van alle vier de buizen na twee uur.
  \item M wordt zoet, M$'$ niet. Formuleer de conclusie, en de precieze
        rol van M$'$.
  \item De blokjes in S krimpen, worden zacht en helder; die in S$'$
        blijven onveranderd. Wat heeft het maagsap gedaan --- en welke
        familie van \omterm{def:g7:digestion-and-nutrients:families}{voedingsstoffen} wordt er voortgebracht?
  \item Er is in geen enkele buis gekauwd of gekneed. Wat bewijst het rek
        over waar het echte werk van de \omterm{def:g4:journey-of-food:digestion}{vertering} ligt?
\end{enumerate}

\medskip\noindent\textbf{Deel II --- Het model verbeteren.}
\begin{enumerate}[resume]
  \item Een leerling stelt voor buis M ``voor de netheid'' op
        koelkasttemperatuur te laten draaien. Maak bezwaar, met het
        patroon uit \cref{prop:g6:decomposers-and-soil:speed} toegepast
        op sappen.
  \item Een ander stelt één grote buis voor met beide sappen door elkaar.
        Welk feit over station-voor-station uit
        \cref{prop:g7:digestion-and-nutrients:chemical} zou dat mengsel
        vertroebelen?
  \item De glazen darm heeft geen wand: er wordt niets opgenomen. Welk
        onderdeel zou een volledig model daarna nodig hebben, en welke
        twee eigenschappen moet het materiaal ervan hebben?
  \item Voeg vezels toe --- zemelen --- aan een vijfde buis met alle
        beschikbare sappen. Voorspel haar toestand na twee uur, en haar
        betekenis voor het model.
\end{enumerate}

\medskip\noindent\textbf{Deel III --- Van glas naar darm.}
\begin{enumerate}[resume]
  \item In het lichaam zouden de suikers van M daarna de \omterm{prop:g7:digestion-and-nutrients:absorption}{darmvlokken}
        tegenkomen. Zet hun reis voort in vier haltes.
  \item Het maagsap werkt in sterk zuur; de sappen van de darm hebben
        nodig dat dat zuur geneutraliseerd wordt. Wat zegt dat over ``elk
        sap bij de \omterm{def:g6:exploring-our-environment:conditions}{omstandigheden} van zijn eigen station'' --- en over de
        stations op volgorde doorslikken?
  \item Een farmaceutisch bedrijf verkoopt capsules met sappen voor
        mensen bij wie de eigen stations tekortschieten. Leg met het rek
        uit waarom zo'n behandeling überhaupt mogelijk is.
  \item Vat het practicum samen in één zin met daarin \emph{sappen},
        \emph{\omterm{def:g7:digestion-and-nutrients:families}{voedingsstoffen}} en \emph{controlebuizen}.
\end{enumerate}
\end{problem}
